% !TeX program = pdflatex
\documentclass[11pt]{article}

\usepackage[margin=1in]{geometry}
\usepackage{amsmath,amssymb,bm}
\usepackage{booktabs}
\usepackage{array}
\usepackage{multirow}
\usepackage{graphicx}
\usepackage{float}
\usepackage{microtype}
\usepackage{siunitx}
\usepackage{xcolor}
\usepackage{hyperref}
\usepackage{enumitem}
\usepackage{longtable}
\usepackage{url}
\usepackage{listings}

\hypersetup{
    colorlinks=true,
    linkcolor=blue,
    citecolor=blue,
    urlcolor=blue
}

\setlist[itemize]{leftmargin=1.5em}
\setlist[enumerate]{leftmargin=1.7em}

\newcommand{\WCT}{\mathrm{WCT}}
\newcommand{\CMS}{\mathrm{CMS}}
\newcommand{\dd}{\mathrm{d}}
\newcommand{\lnm}{\ln(m_{\mu\mu}/m_0)}
\newcommand{\chisq}{\chi^2}
\newcommand{\dchisq}{\Delta\chi^2}

\title{\textbf{Cross-Period CMS Open-Data Replication of a Pre-Existing\\
Wave Confinement Theory Log-Periodic Collider Prediction}}

\author{Richard J. Reyes\\
Wave Confinement Theory / WaveLock\\
\texttt{rickyjreyes.github.io}}

\date{August 28, 2026}

\begin{document}

\maketitle

\begin{abstract}
Wave Confinement Theory (WCT) proposed prior to the present CMS analysis that
curvature-induced spectral organization could generate discrete-scale,
log-periodic structure in collider observables.  In particular, a WCT preprint
published on April 23, 2026 proposed a log-periodic deformation of the
semileptonic Wilson coefficient \(C_9(q^2)\) of the form
\[
\delta C_9(q^2)
=
A\cos\!\left[
\omega\ln\!\left(\frac{q^2}{\Lambda^2}\right)+\phi
\right].
\]
The present work does not measure \(C_9\) directly.  Instead, it performs an
independent, model-agnostic search for related log-periodic structure in the
inclusive opposite-sign dimuon invariant-mass spectrum of CMS Run-2 open data.

A development-stage scan over \(2\)--\(120~\mathrm{GeV}\) initially selected
the lower scan boundary, revealing a broad-background degeneracy rather than a
resolved oscillation.  The analysis was therefore tightened by excluding
known resonance windows from both continuum fitting and spectral testing,
requiring frequencies corresponding to approximately two or more cycles
across the analyzed log-mass range, applying the official CMS 2016 muon
certification, and introducing both residual-permutation and parametric
Poisson background-refit nulls.

In the certified Run2016H discovery file, an interior candidate was found at
\[
\omega_m = 7.0258258258,
\]
with amplitude \(A_r=0.7543\), phase \(\phi=-0.1890\), and
\(\Delta\chi^2=75.76\).  The candidate and complete analysis configuration
were then frozen before independent-file testing.  A second Run2016H file
reproduced the frozen frequency with \(A_r=0.9367\),
\(\phi=-0.3060\), and \(\Delta\chi^2=118.91\).  More importantly, a separately
preregistered Run2016G cross-period test reproduced the same frozen frequency
with \(A_r=0.9349\), \(\phi=-0.1568\), and
\(\Delta\chi^2=115.89\).  In both replication samples, zero of 1000 residual
permutations and zero of 500 parametric Poisson background-refit trials
exceeded the observed fixed-frequency statistic, so the corresponding
reported Monte Carlo \(p\)-values are resolution floors rather than precise
tail probabilities.

These results constitute positive cross-period empirical support for the
pre-existing WCT prediction class of collider log-periodicity.  They do not,
by themselves, establish WCT as the unique explanation.  The principal
remaining alternatives include stable CMS trigger/acceptance/reconstruction
structure, incomplete Standard Model or continuum modeling, and generic
discrete-scale mechanisms already known in the literature.  The decisive next
stage is therefore not further exploratory frequency searching, but frozen
tests across background families, trigger-stable regions, independent
channels/experiments, and a WCT-derived dimensionless mapping that predicts
frequency and phase before new data are inspected.
\end{abstract}

\section{Introduction}

Discrete-scale invariance (DSI) is a well-established mathematical mechanism
by which continuous scaling is replaced by invariance under preferred scale
ratios, naturally producing periodic corrections in the logarithm of a scale
variable \cite{Sornette1998}.  Log-periodic behavior has consequently appeared
in many theoretical and empirical contexts.  In high-energy physics, prior
examples include DSI-inspired descriptions of hadron mass spectra
\cite{Tatischeff2011}, log-periodic decorations of LHC transverse-momentum
spectra within complex-\(q\) or Tsallis-type descriptions
\cite{Wilk2014,Wilk2015}, and approximately periodic invariant-mass resonance
towers in clockwork/linear-dilaton models \cite{Giudice2018}.

Wave Confinement Theory (WCT) makes a more specific claim: confinement and
curvature-driven spectral organization can produce discrete-scale structure
as an emergent property of confined wavefields.  A WCT preprint published on
April 23, 2026 applied this idea specifically to the
\(b\to s\ell^+\ell^-\) sector and proposed a log-periodic deformation of
\(C_9(q^2)\) \cite{ReyesC9}.  Thus the general prediction of collider
log-periodicity predates the CMS analysis reported here.

The present analysis asks a deliberately narrower empirical question:
\begin{quote}
Does an inclusive CMS opposite-sign dimuon mass spectrum contain a stable,
resolved log-periodic residual that survives certification, explicit resonance
exclusions, background-refit null tests, a frozen independent-file
replication, and a frozen cross-period replication?
\end{quote}

The answer obtained under the present pipeline is positive.  However, the
interpretation must distinguish three levels:
\begin{enumerate}
    \item detection of reproducible log-periodic structure in the analyzed CMS
          observable;
    \item consistency with a pre-existing WCT prediction class;
    \item unique attribution of that structure to WCT rather than detector,
          Standard Model, statistical, or other DSI mechanisms.
\end{enumerate}
The first two are supported by the current analysis.  The third remains open.

\section{Prior WCT prediction}

\subsection{Broader WCT structure}

A central scalar curvature functional used in WCT is
\begin{equation}
\Theta[\psi]
=
-\frac{\nabla^2\psi}
{\psi+\epsilon e^{-\alpha|\psi|^2}},
\label{eq:theta}
\end{equation}
with a positive energy functional of the schematic form
\begin{equation}
E[\psi]
=
\int
\left(
|\nabla\psi|^2 + |\Theta[\psi]|^2
\right)\,\dd^n x.
\label{eq:energy}
\end{equation}
The WCT program investigates whether nonlinear curvature feedback, finite-band
spectral selection, winding structure, and confinement can organize physical
observables into preferred spectral structures.

The present CMS work does \emph{not} derive
\(\omega_m=7.0258\) from Eqs.~\eqref{eq:theta}--\eqref{eq:energy}.  That is an
important open theoretical step.  What was predicted prior to this CMS test is
the \emph{class} of log-periodic collider structure.

\subsection{April 2026 collider prediction}

The April 2026 WCT preprint proposed
\begin{equation}
\delta C_9(q^2)
=
A\cos\!\left[
\omega\ln\!\left(\frac{q^2}{\Lambda^2}\right)+\phi
\right],
\label{eq:c9}
\end{equation}
motivated by curvature-induced spectral structure in a confined wavefield
\cite{ReyesC9}.  In conceptual form, the proposed chain was
\begin{equation}
\text{curvature feedback}
\longrightarrow
\text{spectral band selection}
\longrightarrow
\text{discrete-scale structure}
\longrightarrow
\text{log-periodic modulation}.
\label{eq:wct_chain}
\end{equation}

Equation~\eqref{eq:c9} was written for the
\(b\to s\ell^+\ell^-\) phenomenology associated with
\(B^0\to K^{*0}\mu^+\mu^-\).  The inclusive CMS DoubleMuon spectrum analyzed
below is therefore \textbf{not a direct measurement of \(C_9\)} and should not
be described as such.  The CMS analysis instead probes whether a related
log-periodic residual occurs in a distinct collider dimuon observable.

\subsection{Coordinate convention}

The CMS analysis uses
\begin{equation}
x_m=\ln\!\left(\frac{m_{\mu\mu}}{m_0}\right),
\end{equation}
whereas the WCT \(C_9\) construction and the associated LHCb analysis use
\begin{equation}
x_{q^2}=\ln(q^2).
\end{equation}
Since \(q^2=m_{\mu\mu}^2\),
\begin{equation}
\ln q^2 = 2\ln m_{\mu\mu},
\end{equation}
and therefore frequencies expressed in the two conventions obey
\begin{equation}
\boxed{
\omega_m=2k_{q^2},
\qquad
k_{q^2}=\frac{\omega_m}{2}.
}
\label{eq:freqconvert}
\end{equation}
The frozen CMS value
\(\omega_m=7.0258258258\) corresponds to
\begin{equation}
k_{q^2}=3.5129129129.
\end{equation}
Raw numerical frequencies from different log-coordinate definitions must not
be compared without this conversion.

\section{Prior art and competing mechanisms}

The existence of log-periodic behavior is not unique to WCT.  The relevant
question is whether a competing mechanism predicts the same observable,
coordinate dependence, frequency/phase stability, and cross-domain behavior.

\subsection{Generic discrete-scale invariance}

DSI generically yields complex scaling exponents and log-periodic corrections
to scaling \cite{Sornette1998}.  Therefore observation of a log-periodic
residual alone cannot establish WCT.  Generic DSI is best regarded as a
mathematical class containing multiple possible physical realizations.

\subsection{Hadron-spectroscopy DSI}

Tatischeff applied fractal/DSI methods to hadron masses and widths, reporting
log-periodic corrections in hadron spectroscopy \cite{Tatischeff2011}.  This
is relevant high-energy prior art, but it is not a prediction of the present
inclusive CMS dimuon continuum residual and does not use the WCT confinement
mechanism.

\subsection{Complex-\(q\)/Tsallis log-periodicity}

Wilk and W{\l}odarczyk reported log-periodic oscillatory decorations in
high-\(p_T\) spectra and discussed descriptions based on a complex
nonextensivity parameter or a log-periodically varying scale
\cite{Wilk2014,Wilk2015}.  This is an important collider-specific alternative
to any broad claim that WCT uniquely introduced collider log-periodicity.
However, the primary observable and physical interpretation differ from the
present dimuon invariant-mass residual.

\subsection{Clockwork/linear-dilaton resonance towers}

Clockwork/linear-dilaton phenomenology predicts a tower of approximately
periodic resonant peaks and motivates Fourier-space searches in dilepton and
diphoton invariant-mass spectra \cite{Giudice2018}.  This is perhaps the
closest prior collider model in terms of a structured dilepton mass spectrum.
Nevertheless, a resonance comb approximately periodic in mass is not the same
functional hypothesis as
\begin{equation}
r(m)\propto
\cos\!\left[\omega\ln(m/m_0)-\phi\right].
\end{equation}
The two can and should be compared quantitatively rather than conflated.

\subsection{Standard Model and detector/systematic structure}

For the present CMS result, the most serious competing explanation is not
necessarily another beyond-Standard-Model theory.  Stable structures can arise
from
\begin{itemize}
    \item trigger thresholds and changing trigger efficiencies;
    \item reconstruction and identification efficiencies;
    \item detector acceptance versus \(m_{\mu\mu}\);
    \item residual resonance tails or heavy-flavor structure;
    \item continuum background misspecification;
    \item correlated residuals not represented by a Poisson null;
    \item run-dependent or run-stable calibration effects.
\end{itemize}
Because Run2016G and Run2016H use closely related CMS reconstruction and trigger
infrastructure, cross-period replication is stronger than a same-file
fluctuation test but does not eliminate CMS-specific systematic explanations.

\section{CMS open data}

\subsection{Run2016H}

The discovery and first frozen replication use the CMS
\texttt{DoubleMuon} Run2016H Ultra Legacy NanoAOD dataset
\cite{CMS2016H}:
\begin{quote}
\texttt{/DoubleMuon/Run2016H-UL2016\_MiniAODv2\_NanoAODv9-v1/NANOAOD}.
\end{quote}
The CERN Open Data record reports 48,912,812 events in 28 files from run
numbers 281613--284044 \cite{CMS2016H}.

The two analyzed files are
\begin{align}
&\texttt{127C2975-1B1C-A046-AABF-62B77E757A86.root},
\\
&\texttt{183BFB78-7B5E-734F-BBF5-174A73020F89.root}.
\end{align}

\subsection{Run2016G}

The cross-period replication uses the CMS
\texttt{DoubleMuon} Run2016G Ultra Legacy NanoAOD dataset
\cite{CMS2016G}:
\begin{quote}
\texttt{/DoubleMuon/Run2016G-UL2016\_MiniAODv2\_NanoAODv9-v2/NANOAOD}.
\end{quote}
The CERN Open Data record reports 45,235,604 events in 29 files from run
numbers 278820--280385 \cite{CMS2016G}.

The preregistered first file used for the cross-period test is
\begin{quote}
\texttt{05DD095C-F6C3-9A4F-9FB3-348A5A6403D5.root}.
\end{quote}

\subsection{Muon certification}

All primary certified analyses apply
\begin{quote}
\texttt{Cert\_271036-284044\_13TeV\_Legacy2016\_Collisions16\_JSON\_MuonPhys.txt},
\end{quote}
the official CMS list of luminosity sections valid for analyses requiring
valid muons \cite{CMSMuonJSON}.  The certification record covers both
Run2016G and Run2016H.

\section{Event reconstruction and selection}

NanoAOD files are read with \texttt{uproot} and \texttt{awkward}.  Muon
four-vectors are reconstructed from the standard NanoAOD quantities.  The
selection used throughout the frozen analysis is
\begin{align}
p_T^\mu &\ge 4~\mathrm{GeV},\\
|\eta^\mu| &\le 2.4,\\
\text{ID} &= \text{tight}.
\end{align}
Opposite-sign muon pairs are formed, and their invariant mass is
\begin{equation}
m_{\mu\mu}^2
=
(E_{\mu^+}+E_{\mu^-})^2
-
\left|
\mathbf{p}_{\mu^+}+\mathbf{p}_{\mu^-}
\right|^2.
\end{equation}

The analysis interval is
\begin{equation}
2~\mathrm{GeV}
\le m_{\mu\mu}\le
120~\mathrm{GeV},
\end{equation}
divided into 350 logarithmically spaced bins.

\subsection{Explicit resonance masks}

The following mass windows are excluded from both the continuum fit and the
spectral analysis:
\begin{align}
2.90 &\le m_{\mu\mu}\le 3.30~\mathrm{GeV},\\
3.55 &\le m_{\mu\mu}\le 3.85~\mathrm{GeV},\\
8.50 &\le m_{\mu\mu}\le 11.50~\mathrm{GeV},\\
80.0 &\le m_{\mu\mu}\le 100.0~\mathrm{GeV}.
\end{align}
These remove the dominant \(J/\psi\), \(\psi(2S)\), bottomonium-family, and
\(Z\)-region structures from the tested continuum.

This exclusion is essential.  Omitting a resonance from a background fit
while retaining it in the residual spectral scan would deliberately generate
a large structured residual.

\section{Continuum model}

Let \(N_i\) denote the observed count in mass bin \(i\) with center \(m_i\).
The smooth background is estimated by fitting
\begin{equation}
y_i=\ln(N_i+0.5)
\end{equation}
as a Chebyshev polynomial in
\begin{equation}
x_i=\ln m_i.
\end{equation}
The frozen polynomial degree is seven.

The fit is iteratively robustified.  At each iteration,
\begin{equation}
r_i
=
\frac{N_i-B_i}{\sqrt{\max(B_i,1)}}
\label{eq:pearson}
\end{equation}
is computed, and bins beyond the declared clipping threshold are excluded
from subsequent continuum fitting.  Explicit resonance masks remain excluded
throughout.

Only bins with a model expectation above the declared minimum count are
retained for spectral testing.

\section{Log-periodic spectral test}

The residual model is
\begin{equation}
r(x)
=
c
+
a\cos(\omega x)
+
b\sin(\omega x),
\label{eq:linmodel}
\end{equation}
equivalently
\begin{equation}
r(x)
=
c
+
A_r\cos(\omega x-\phi),
\end{equation}
where
\begin{align}
A_r &= \sqrt{a^2+b^2},\\
\phi &= \operatorname{atan2}(b,a).
\end{align}

The scan statistic is the improvement relative to an intercept-only residual
model:
\begin{equation}
\Delta\chi^2(\omega)
=
\chi^2_0-\chi^2_\omega.
\end{equation}

The exploratory scan is
\begin{equation}
3.1\le\omega\le80
\end{equation}
with 1000 linearly spaced trial frequencies.

\subsection{Spectral-resolution condition}

For a log-coordinate span
\begin{equation}
\Delta x=x_{\max}-x_{\min},
\end{equation}
the number of cycles represented by a trial frequency is
\begin{equation}
N_{\mathrm{cycles}}
=
\frac{\omega\Delta x}{2\pi}.
\label{eq:ncycles}
\end{equation}
Across the approximately
\(\Delta x\simeq4.08\) analyzed span, \(\omega\simeq3.1\) corresponds to
approximately two cycles.  Frequencies below this scale were excluded from the
final exploratory search because sub-cycle behavior is strongly degenerate
with broad continuum curvature.

\section{Null distributions}

Two Monte Carlo null procedures are used.

\subsection{Residual-permutation null}

The analyzed residuals are permuted across the fixed log-mass coordinates.
For every permutation the complete frequency scan is repeated and its maximum
\(\Delta\chi^2\) is recorded.  This controls the look-elsewhere effect over the
declared \(\omega\)-scan under the assumption that fitted residuals are
exchangeable.

For a frozen-frequency replication, the statistic at the predeclared
frequency is also evaluated under each permutation.

The Monte Carlo estimate uses the add-one correction
\begin{equation}
\hat p
=
\frac{N_{\ge}+1}{N_{\mathrm{MC}}+1}.
\label{eq:addone}
\end{equation}

\subsection{Parametric Poisson background-refit null}

The second null propagates Poisson counting noise and continuum refitting.
For each pseudo-experiment:
\begin{enumerate}
    \item counts are drawn as
    \begin{equation}
    N_i^\ast\sim\operatorname{Poisson}(B_i);
    \end{equation}
    \item the degree-seven robust continuum is refit to \(N_i^\ast\);
    \item new Pearson residuals are computed;
    \item the exploratory scan and/or frozen-frequency statistic is evaluated.
\end{enumerate}

This is stronger than simply shuffling residual bins because the continuum is
re-estimated in every pseudo-experiment.  It still does \emph{not} simulate
trigger, reconstruction, calibration, acceptance, correlated detector effects,
or an alternative Standard Model continuum family.

\section{Analysis development and rejection of the initial boundary solution}

The first 100,000-event development run used a broad \(2\)--\(120~\mathrm{GeV}\)
fit without the final resonance exclusions and scanned
\(0.5\le\omega\le80\).  It selected
\begin{equation}
\omega_{\mathrm{best}}=0.5,
\end{equation}
exactly the lower scan boundary, with an unphysical residual amplitude of
order \(10^2\).

The full log span is
\begin{equation}
\ln(120/2)=\ln 60\simeq4.094,
\end{equation}
so \(\omega=0.5\) represents only
\begin{equation}
N_{\mathrm{cycles}}
\simeq
\frac{0.5(4.094)}{2\pi}
\simeq0.326
\end{equation}
cycles.  The candidate was therefore interpreted as broad background
misspecification rather than a resolved oscillatory signal.

This negative development-stage finding motivated explicit safeguards:
\begin{itemize}
    \item resonance windows excluded from both fitting and testing;
    \item an exploratory lower bound near the two-cycle resolution scale;
    \item explicit boundary-frequency warnings;
    \item residual RMS and maximum-residual diagnostics;
    \item official CMS muon certification;
    \item background-refit parametric nulls.
\end{itemize}

The initial \(\omega=0.5\) solution is not used as evidence for WCT.

\section{Certified Run2016H discovery}

With the final masks and official muon certification applied to the first
Run2016H file, 100,000 events were read and 50,951 survived the certification
within that input segment.  The analysis contained 21,930 dimuon masses in the
declared mass interval and 287 analyzed bins.

The exploratory scan selected
\begin{equation}
\boxed{
\omega_m=7.0258258258
}
\label{eq:frozenomega}
\end{equation}
with
\begin{align}
A_r &= 0.7542594,\\
\phi &= -0.1889538,\\
\Delta\chi^2 &= 75.7616.
\end{align}
The candidate spans
\begin{equation}
N_{\mathrm{cycles}}=4.5652,
\end{equation}
and is not on a scan boundary.

The residual diagnostics were
\begin{align}
\mathrm{RMS}(r)&=1.2865,\\
\max|r|&=6.3986.
\end{align}
Thus the residuals are substantially better behaved than in the original
development scan but are not perfectly unit-normal.

Zero of 1000 residual-permutation trials exceeded the observed global scan
maximum, and zero of 500 Poisson background-refit trials exceeded it.  Using
Eq.~\eqref{eq:addone}, the reported Monte Carlo floors are
\begin{align}
\hat p_{\mathrm{perm}} &= \frac{1}{1001}
=9.9900\times10^{-4},\\
\hat p_{\mathrm{refit}} &= \frac{1}{501}
=1.9960\times10^{-3}.
\end{align}
These values mean zero exceedances at the available Monte Carlo resolution;
they are not precision measurements of the true tail probability.

\section{Preregistration of the CMS replication}

Before inspecting a second CMS file, the following were frozen:
\begin{itemize}
    \item \(\omega_m=7.0258258258\);
    \item mass range \(2\)--\(120~\mathrm{GeV}\);
    \item 350 logarithmic bins;
    \item \(p_T^\mu\ge4~\mathrm{GeV}\);
    \item \(|\eta^\mu|\le2.4\);
    \item tight muon ID;
    \item degree-seven continuum;
    \item the four explicit mass exclusions;
    \item 100,000-event read cap;
    \item 1000 residual permutations;
    \item 500 parametric background-refit trials;
    \item random seed 20260827.
\end{itemize}

The exploratory best frequency in a replication file is secondary.  The
primary replication statistic is the amplitude and \(\Delta\chi^2\) obtained
at the frozen frequency in Eq.~\eqref{eq:frozenomega}.

\section{Frozen Run2016H independent-file replication}

The second Run2016H ROOT file was not used to determine the frozen frequency.
After applying the same configuration, 100,000 events were read, 90,868
survived certification, and 39,230 masses fell within the full mass interval.

The exploratory scan of this replication file independently selected
\begin{equation}
\omega_{\mathrm{explore,H2}}=6.563964,
\end{equation}
but the preregistered primary result is the fit at
\(\omega_m=7.0258258258\):
\begin{align}
A_r &= 0.9367121,\\
\phi &= -0.3059911,\\
\Delta\chi^2 &= 118.9148.
\end{align}

Again,
\begin{align}
N_{\ge,\mathrm{perm}}&=0/1000,\\
N_{\ge,\mathrm{refit}}&=0/500.
\end{align}

The phase shift relative to discovery is
\begin{equation}
|\Delta\phi_{\mathrm{H2-H1}}|
=
0.11704~\mathrm{rad}
\simeq6.71^\circ.
\end{equation}

The principal warning is that the absolute continuum residual diagnostics are
worse:
\begin{align}
\mathrm{RMS}(r)&=1.8187,\\
\max|r|&=12.4685.
\end{align}
Consequently, this sample strongly rejects the implemented Poisson smooth-null
model but simultaneously demonstrates that the smooth background is not a
complete statistical description of the data.

\section{Preregistered Run2016G cross-period replication}

Before inspecting Run2016G data, the same fixed frequency and analysis
configuration were preregistered for a cross-period test.

In the first Run2016G file, 100,000 events were read and all 100,000 were
within accepted muon-certified luminosity sections for the processed segment.
There were 39,633 masses in the declared \(2\)--\(120~\mathrm{GeV}\) interval.

The secondary exploratory scan selected
\begin{equation}
\omega_{\mathrm{explore,G}}=7.410711,
\end{equation}
whereas the primary fixed-frequency result at
\(\omega_m=7.0258258258\) was
\begin{align}
A_r &= 0.9348797,\\
\phi &= -0.1567923,\\
\Delta\chi^2 &= 115.8921.
\end{align}

For the frozen statistic,
\begin{align}
N_{\ge,\mathrm{perm}}&=0/1000,\\
N_{\ge,\mathrm{refit}}&=0/500.
\end{align}

The frozen phase differs from the original discovery phase by only
\begin{equation}
|\Delta\phi_{\mathrm{G-H1}}|
=
0.03216~\mathrm{rad}
\simeq1.84^\circ.
\end{equation}
The amplitude differs by approximately \(24\%\) from the discovery value and
is almost identical to the independent Run2016H replication amplitude.

The Run2016G residual diagnostics are
\begin{align}
\mathrm{RMS}(r)&=1.7240,\\
\max|r|&=8.4725.
\end{align}
As in the second Run2016H file, these values argue against treating the fitted
continuum as a complete physical null model.

\section{Summary of results}

\begin{table}[H]
\centering
\small
\caption{Certified discovery and frozen replication results.  ``H1'' denotes
the Run2016H discovery file, ``H2'' the independent Run2016H replication file,
and ``G'' the preregistered Run2016G cross-period file.}
\label{tab:main}
\begin{tabular}{lrrr}
\toprule
Quantity & H1 discovery & H2 frozen replication & G frozen replication\\
\midrule
Events read & 100000 & 100000 & 100000\\
Events after JSON & 50951 & 90868 & 100000\\
Masses in range & 21930 & 39230 & 39633\\
Analysis bins & 287 & 287 & 287\\
Frozen \(\omega_m\) & 7.025826 & 7.025826 & 7.025826\\
Amplitude \(A_r\) & 0.754259 & 0.936712 & 0.934880\\
Phase \(\phi\) [rad] & -0.188954 & -0.305991 & -0.156792\\
\(\Delta\chi^2\) & 75.7616 & 118.9148 & 115.8921\\
Residual RMS & 1.2865 & 1.8187 & 1.7240\\
Maximum \(|r|\) & 6.3986 & 12.4685 & 8.4725\\
Permutation exceedances & 0/1000 & 0/1000 & 0/1000\\
Refit-bootstrap exceedances & 0/500 & 0/500 & 0/500\\
\bottomrule
\end{tabular}
\end{table}

The three frozen-frequency amplitudes have mean
\begin{equation}
\overline{A}_r \simeq 0.8753
\end{equation}
with a sample coefficient of variation of approximately \(12\%\).
No combined significance is assigned because the samples share analysis
choices and, in particular, the CMS periods share detector/reconstruction
infrastructure.

\begin{table}[H]
\centering
\small
\caption{Exploratory best frequencies.  These are secondary in the two
replication datasets because the primary test is fixed at the preregistered
\(\omega_m=7.0258258258\).}
\begin{tabular}{lrrr}
\toprule
Sample & Exploratory best \(\omega\) & Best \(\Delta\chi^2\) & Best amplitude\\
\midrule
Run2016H discovery & 7.025826 & 75.7616 & 0.7543\\
Run2016H file 2 & 6.563964 & 128.5539 & 0.9719\\
Run2016G file 1 & 7.410711 & 144.1300 & 1.0316\\
\bottomrule
\end{tabular}
\end{table}
\section{Figures}

The following figures summarize the certified Run2016H discovery,
the independent Run2016H replication, and the preregistered Run2016G
cross-period replication.  For replication samples, the frequency scan is
preferred over the automatically generated
\texttt{best\_log\_periodic\_fit.png}, because the latter displays the
exploratory best-fit frequency rather than the preregistered frozen frequency.

\begin{figure}[H]
\centering
\IfFileExists{results/dimuon_certified_bootstrap500/omega_scan.png}{
\includegraphics[width=0.92\linewidth]
{results/dimuon_certified_bootstrap500/omega_scan.png}
}{
\fbox{\parbox[c][2.0in][c]{0.84\linewidth}{\centering
Place
\texttt{results/dimuon\_certified\_bootstrap500/omega\_scan.png}
here.}}
}
\caption{
Certified Run2016H discovery frequency scan.
The dominant resolved interior candidate occurs at
$\omega_m=7.0258258258$, with
$\Delta\chi^2=75.76$ and an oscillation spanning approximately
$4.57$ cycles across the analyzed logarithmic mass interval.
This value was subsequently frozen before the independent CMS
replication tests.
}
\label{fig:run2016h_discovery_scan}
\end{figure}


\begin{figure}[H]
\centering
\IfFileExists{results/dimuon_replication_file2_frozen/omega_scan.png}{
\includegraphics[width=0.92\linewidth]
{results/dimuon_replication_file2_frozen/omega_scan.png}
}{
\fbox{\parbox[c][2.0in][c]{0.84\linewidth}{\centering
Place
\texttt{results/dimuon\_replication\_file2\_frozen/omega\_scan.png}
here.}}
}
\caption{
Independent Run2016H replication frequency scan.
The dotted vertical line marks the preregistered frozen frequency
$\omega_m=7.0258258258$ obtained from the discovery sample.
Although the secondary exploratory maximum occurs at
$\omega=6.5640$, the frozen frequency itself remains strongly supported,
with $\Delta\chi^2=118.91$, amplitude $A_r=0.937$, and zero exceedances
in both 1000 residual-permutation trials and 500 parametric
background-refit trials.
}
\label{fig:run2016h_replication_scan}
\end{figure}


\begin{figure}[H]
\centering
\IfFileExists{results/dimuon_run2016g_frozen/omega_scan.png}{
\includegraphics[width=0.92\linewidth]
{results/dimuon_run2016g_frozen/omega_scan.png}
}{
\fbox{\parbox[c][2.0in][c]{0.84\linewidth}{\centering
Place
\texttt{results/dimuon\_run2016g\_frozen/omega\_scan.png}
here.}}
}
\caption{
Preregistered Run2016G cross-period replication frequency scan.
The dotted vertical line marks the frozen prediction
$\omega_m=7.0258258258$, fixed before inspection of the Run2016G
sample.  The independently selected exploratory maximum occurs at
$\omega=7.4107$, while the frozen frequency remains within the dominant
spectral structure and gives
$\Delta\chi^2=115.89$ with amplitude $A_r=0.935$.
Zero of 1000 residual permutations and zero of 500 parametric
background-refit trials exceeded the observed frozen-frequency statistic.
}
\label{fig:run2016g_scan}
\end{figure}


\begin{figure}[H]
\centering
\IfFileExists{results/dimuon_run2016g_frozen/spectrum_background.png}{
\includegraphics[width=0.92\linewidth]
{results/dimuon_run2016g_frozen/spectrum_background.png}
}{
\fbox{\parbox[c][2.0in][c]{0.84\linewidth}{\centering
Place
\texttt{results/dimuon\_run2016g\_frozen/spectrum\_background.png}
here.}}
}
\caption{
Certified Run2016G opposite-sign dimuon invariant-mass spectrum and
the robust degree-seven smooth background model used to construct the
Pearson residuals.  The dominant $J/\psi$, $\psi(2S)$, bottomonium,
and $Z$ mass windows are excluded from the subsequent spectral analysis.
}
\label{fig:run2016g_spectrum}
\end{figure}


\begin{figure}[H]
\centering
\IfFileExists{results/dimuon_run2016g_frozen/residuals.png}{
\includegraphics[width=0.92\linewidth]
{results/dimuon_run2016g_frozen/residuals.png}
}{
\fbox{\parbox[c][2.0in][c]{0.84\linewidth}{\centering
Place
\texttt{results/dimuon\_run2016g\_frozen/residuals.png}
here.}}
}
\caption{
Pearson residuals of the certified Run2016G dimuon spectrum after
the smooth-background fit and declared mass exclusions.
The residual RMS is $1.724$ and the maximum absolute residual is
$8.473$.  The remaining structured variation demonstrates that the
degree-seven continuum is not a complete physical null model and
motivates the additional systematic controls discussed in the text.
}
\label{fig:run2016g_residuals}
\end{figure}


\begin{figure}[H]
\centering
\IfFileExists{results/dimuon_run2016g_frozen/global_null_distribution.png}{
\includegraphics[width=0.88\linewidth]
{results/dimuon_run2016g_frozen/global_null_distribution.png}
}{
\fbox{\parbox[c][2.0in][c]{0.84\linewidth}{\centering
Place
\texttt{results/dimuon\_run2016g\_frozen/global\_null\_distribution.png}
here.}}
}
\caption{
Residual-permutation distribution of the maximum exploratory
$\Delta\chi^2$ obtained across the declared frequency scan in Run2016G.
No permutation exceeded the observed scan maximum.
This figure represents the global exploratory null distribution;
the preregistered frozen-frequency significance is separately obtained
from the fixed-frequency permutation and parametric-refit statistics.
}
\label{fig:run2016g_global_null}
\end{figure}

\section{Interpretation for WCT}

\subsection{What the data support}

The CMS results support the following statement:
\begin{quote}
A log-periodic residual class proposed by WCT prior to the present CMS analysis
is observed in a distinct inclusive CMS dimuon observable, and a
CMS-discovered numerical candidate within that class subsequently reproduces
at a preregistered fixed frequency in both an independent Run2016H file and a
separately preregistered Run2016G cross-period sample.
\end{quote}

This is stronger than a single exploratory pattern match for four reasons:
\begin{enumerate}
    \item the broad low-frequency boundary artifact was explicitly rejected;
    \item the surviving candidate is an interior spectral feature spanning
          more than four cycles;
    \item the second Run2016H test used a frozen frequency;
    \item the Run2016G test was frozen before that period was inspected.
\end{enumerate}

The result is therefore properly described as
\begin{quote}
\textbf{positive cross-period empirical support for the pre-existing WCT
prediction class of collider log-periodicity.}
\end{quote}

\subsection{What the data do not yet establish}

The present result does not establish:
\begin{itemize}
    \item that WCT is the unique origin of the CMS structure;
    \item that the CMS inclusive dimuon residual is a measurement of the
          \(C_9(q^2)\) deformation in Eq.~\eqref{eq:c9};
    \item that \(\omega_m=7.0258\) was derived from WCT field equations before
          the CMS data were analyzed;
    \item that the observed modulation survives a complete detector,
          trigger-efficiency, acceptance, and Standard Model treatment;
    \item that the same dimensionless WCT frequency has already been
          demonstrated across CMS, LHCb, gravitational-wave, photodiode, and
          atomic datasets.
\end{itemize}

These distinctions are not rhetorical qualifications; they define the next
falsifiable tests.

\section{Why the CMS result is not yet ``case closed''}

A repeated fixed-frequency signal can arise from either a reproducible physical
effect or a reproducible systematic effect.  Run2016G and Run2016H contain
different collisions, and the cross-period replication therefore strongly
weakens a one-off statistical-fluctuation explanation.  However, both periods
remain CMS 2016 data processed through closely related trigger, detector, and
reconstruction systems.

The elevated residual RMS values,
\begin{equation}
1.29,\quad1.82,\quad1.72,
\end{equation}
also show that the degree-seven smooth continuum does not fully whiten the
spectrum.  The refit-bootstrap result establishes that the observed fixed
frequency is not reproduced by \emph{Poisson fluctuations around that fitted
continuum}.  It does not establish that the fitted continuum is the correct
complete null.

The strongest remaining mundane hypothesis is therefore
\begin{equation}
\boxed{
\text{stable mass-dependent CMS/SM structure not represented by the
current smooth null}.
}
\end{equation}

\section{Specific systematic risks}

\subsection{Trigger and acceptance}

The \texttt{DoubleMuon} dataset is selected by dimuon trigger paths and the
current analysis does not yet enforce a single demonstrably flat trigger
plateau over the entire \(2\)--\(120~\mathrm{GeV}\) interval or apply a
mass-dependent trigger-efficiency correction.  This is a major outstanding
control.  A stable efficiency modulation can reproduce across files and run
periods.

\subsection{Muon efficiency and reconstruction}

Tight-ID, tracking, and reconstruction efficiencies can vary with muon
kinematics.  Since \(m_{\mu\mu}\) is correlated with the individual muon
\(p_T\) and \(\eta\), efficiency structure can project into mass space.

\subsection{Continuum model dependence}

A single Chebyshev degree is not sufficient for a final physical inference.
The result must survive a predeclared family of reasonable backgrounds without
selecting the background that maximizes the WCT statistic.

\subsection{Residual resonance and heavy-flavor structure}

Explicit resonance windows remove the dominant peaks, but broad tails,
interference, heavy-flavor components, and transitions between physical
production regimes remain possible.

\subsection{Correlated residuals}

Neither independent residual permutations nor simple Poisson pseudo-counts
model arbitrary bin-to-bin detector/systematic covariance.  A final result
requires structured nulls or simulation/control samples that encode such
correlations.

\section{Predeclared robustness matrix}

A natural next test is a robustness matrix that is frozen before inspection.
For example:
\begin{align}
\text{background degree} &\in \{5,6,7,8,9\},\\
N_{\mathrm{bins}} &\in \{250,300,350,400,450\},\\
\text{binning} &\in \{\text{logarithmic},\text{linear control}\},
\end{align}
with the fixed frequency
\begin{equation}
\omega_m=7.0258258258
\end{equation}
used as the primary statistic in every cell.

The purpose is not to demand that every implementation give identical
significance.  The relevant questions are:
\begin{enumerate}
    \item Does the fitted amplitude retain a consistent sign and order of
          magnitude?
    \item Does the phase remain stable?
    \item Does the fixed-frequency statistic remain unusual under refitted
          nulls?
    \item Is the result concentrated in a narrow analysis choice?
\end{enumerate}

A family-wise null over the entire frozen robustness matrix would provide a
much stronger statement than selecting the best-performing configuration.

\section{Critical external replication tests}

The most informative future tests are those that separate WCT from
CMS-specific structure.

\subsection{Different CMS channel}

A dielectron, photon, or other independent observable could be transformed
into a WCT-prescribed dimensionless coordinate and tested at a frequency
frozen before inspection.  Different trigger and reconstruction systems would
substantially reduce shared-systematic explanations.

\subsection{Independent experiment}

An ATLAS or other independent collider dataset would be substantially more
decisive because CMS-specific efficiency, reconstruction, and calibration
features would no longer be shared.

\subsection{Theory-derived frequency mapping}

The most important theoretical milestone is to derive
\begin{equation}
\boxed{
\text{WCT field parameters}
\longrightarrow
\text{dimensionless scale variable}
\longrightarrow
\omega,\phi,A
}
\end{equation}
before opening a new dataset.

This would convert WCT from predicting a broad phenomenological class
(log-periodicity) to predicting a numerical cross-domain invariant.  Such a
test would distinguish WCT much more sharply from generic DSI.

\section{Relation to the broader multi-domain WCT program}

WCT has also been tested against gravitational-wave catalogs, laboratory
photodiode data, atomic/NIST-related data, and LHC/LHCb observables.  Those
datasets should not be combined merely because each contains a spectral or
log-periodic feature.  A scientifically meaningful cross-domain claim
requires:
\begin{enumerate}
    \item a theory-defined dimensionless coordinate in every domain;
    \item correct conversion between \(\ln m\), \(\ln q^2\), time, frequency,
          energy, or other coordinates;
    \item frozen predictions before each new dataset is inspected;
    \item explicit negative results and failed controls;
    \item a joint null model that accounts for the number of datasets,
          observables, transforms, and analysis families searched.
\end{enumerate}

If WCT derives a single invariant mapping and that mapping predicts the
frequency and phase observed in multiple independently instrumented systems,
the evidentiary status would be qualitatively stronger than independent
post-hoc spectral matches.

\section{Reproducibility}

The analysis implementation is maintained in the
\texttt{wct-cms} repository.  The repository records:
\begin{itemize}
    \item the discovery analysis;
    \item the frozen candidate document;
    \item the independent Run2016H replication result;
    \item the preregistered Run2016G replication;
    \item regression tests for Awkward-array selection;
    \item vectorized spectral scanning;
    \item residual-permutation nulls;
    \item parametric Poisson background-refit nulls.
\end{itemize}

At the time of the reported runs, the local regression suite returned
\begin{quote}
\texttt{8 passed}.
\end{quote}

\subsection{Discovery-file certified run}

\begin{lstlisting}[basicstyle=\ttfamily\small,breaklines=true]
cms-wct \
  --input data/files.txt \
  --golden-json \
  data/cert/14221/Cert_271036-284044_13TeV_Legacy2016_Collisions16_JSON_MuonPhys.txt \
  --output-dir results/dimuon_certified_bootstrap500 \
  --mass-min 2 \
  --mass-max 120 \
  --bins 350 \
  --log-bins \
  --muon-pt-min 4 \
  --muon-eta-max 2.4 \
  --tight-id \
  --fit-degree 7 \
  --mask-window 2.9:3.3 \
  --mask-window 3.55:3.85 \
  --mask-window 8.5:11.5 \
  --mask-window 80:100 \
  --omega-min 3.1 \
  --omega-max 80 \
  --omega-steps 1000 \
  --max-events 100000 \
  --permutations 1000 \
  --parametric-bootstrap 500 \
  --seed 20260827
\end{lstlisting}

\subsection{Frozen Run2016H replication}

\begin{lstlisting}[basicstyle=\ttfamily\small,breaklines=true]
cms-wct \
  --input data/files_replication.txt \
  --golden-json \
  data/cert/14221/Cert_271036-284044_13TeV_Legacy2016_Collisions16_JSON_MuonPhys.txt \
  --output-dir results/dimuon_replication_file2_frozen \
  --mass-min 2 \
  --mass-max 120 \
  --bins 350 \
  --log-bins \
  --muon-pt-min 4 \
  --muon-eta-max 2.4 \
  --tight-id \
  --fit-degree 7 \
  --mask-window 2.9:3.3 \
  --mask-window 3.55:3.85 \
  --mask-window 8.5:11.5 \
  --mask-window 80:100 \
  --omega-min 3.1 \
  --omega-max 80 \
  --omega-steps 1000 \
  --frozen-omega 7.025825825825827 \
  --max-events 100000 \
  --permutations 1000 \
  --parametric-bootstrap 500 \
  --seed 20260827
\end{lstlisting}

\subsection{Frozen Run2016G replication}

\begin{lstlisting}[basicstyle=\ttfamily\small,breaklines=true]
cms-wct \
  --input data/files_run2016g.txt \
  --golden-json \
  data/cert/14221/Cert_271036-284044_13TeV_Legacy2016_Collisions16_JSON_MuonPhys.txt \
  --output-dir results/dimuon_run2016g_frozen \
  --mass-min 2 \
  --mass-max 120 \
  --bins 350 \
  --log-bins \
  --muon-pt-min 4 \
  --muon-eta-max 2.4 \
  --tight-id \
  --fit-degree 7 \
  --mask-window 2.9:3.3 \
  --mask-window 3.55:3.85 \
  --mask-window 8.5:11.5 \
  --mask-window 80:100 \
  --omega-min 3.1 \
  --omega-max 80 \
  --omega-steps 1000 \
  --frozen-omega 7.025825825825827 \
  --max-events 100000 \
  --permutations 1000 \
  --parametric-bootstrap 500 \
  --seed 20260827
\end{lstlisting}

\section{Conclusion}

The CMS analysis reported here began with a failed development-stage
low-frequency boundary solution, which was rejected as a continuum-model
degeneracy.  After explicit resonance masking, a minimum spectral-resolution
criterion, official muon certification, and stronger null procedures, a
resolved Run2016H candidate emerged at
\[
\omega_m=7.0258258258.
\]
That numerical candidate was frozen before an independent CMS file was
inspected and reproduced at the fixed frequency in a second Run2016H file.
The same frequency was then preregistered before a Run2016G test and again
produced a large fixed-frequency statistic, closely reproducing the original
phase:
\[
|\Delta\phi_{\mathrm{G-H1}}|\simeq1.84^\circ.
\]
In both replication samples, zero of 1000 residual permutations and zero of
500 Poisson background-refit pseudo-experiments exceeded the observed
fixed-frequency statistic.

Because WCT had publicly proposed collider log-periodicity before this CMS
analysis, these results are \textbf{positive empirical support for a
pre-existing WCT prediction class}.  They are not yet a unique confirmation
of WCT.  Generic DSI, collider log-periodicity models, and especially stable
CMS/Standard-Model mass-dependent structure remain viable alternatives.
The elevated residual RMS in the replication samples makes that caveat
quantitatively important.

The strongest next test is consequently not another unrestricted spectral
search.  It is a frozen discriminating program: trigger-stable selections,
predeclared background-family robustness, independent detector channels,
independent experiments, and a first-principles WCT mapping that predicts a
dimensionless frequency and phase before new observations are opened.

\begin{thebibliography}{99}

\bibitem{ReyesC9}
R.~J. Reyes,
\newblock ``A Curvature-Induced Log-Periodic Deformation of \(C_9(q^2)\):
Wave Confinement Theory and the LHCb
\(B^0\to K^{*0}\mu^+\mu^-\) Anomaly,''
\newblock Zenodo, Version 1, published April 23, 2026,
\newblock DOI: \href{https://doi.org/10.5281/zenodo.19705254}
{10.5281/zenodo.19705254}.

\bibitem{CMS2016H}
CMS Collaboration,
\newblock ``DoubleMuon primary dataset in NANOAOD format from RunH of 2016
(\texttt{/DoubleMuon/Run2016H-UL2016\_MiniAODv2\_NanoAODv9-v1/NANOAOD}),''
\newblock CERN Open Data Portal, 2024,
\newblock DOI:
\href{https://doi.org/10.7483/OPENDATA.CMS.UZD7.Z50M}
{10.7483/OPENDATA.CMS.UZD7.Z50M}.

\bibitem{CMS2016G}
CMS Collaboration,
\newblock ``DoubleMuon primary dataset in NANOAOD format from RunG of 2016
(\texttt{/DoubleMuon/Run2016G-UL2016\_MiniAODv2\_NanoAODv9-v2/NANOAOD}),''
\newblock CERN Open Data Portal, 2024,
\newblock DOI:
\href{https://doi.org/10.7483/OPENDATA.CMS.ZQS3.LGLP}
{10.7483/OPENDATA.CMS.ZQS3.LGLP}.

\bibitem{CMSMuonJSON}
CMS Collaboration,
\newblock ``CMS list of validated runs
\texttt{Cert\_271036-284044\_13TeV\_Legacy2016\_Collisions16\_JSON\_MuonPhys.txt},''
\newblock CERN Open Data Portal, 2023,
\newblock \url{https://opendata.cern.ch/record/cms-14221}.

\bibitem{Sornette1998}
D.~Sornette,
\newblock ``Discrete-scale invariance and complex dimensions,''
\newblock \emph{Physics Reports}, vol.~297, no.~5, pp.~239--270, 1998,
\newblock DOI:
\href{https://doi.org/10.1016/S0370-1573(97)00076-8}
{10.1016/S0370-1573(97)00076-8}.

\bibitem{Tatischeff2011}
B.~Tatischeff,
\newblock ``Fractal properties applied to hadron spectroscopy,''
\newblock arXiv:1105.2034, 2011,
\newblock \url{https://arxiv.org/abs/1105.2034}.

\bibitem{Wilk2014}
G.~Wilk and Z.~W{\l}odarczyk,
\newblock ``Log-periodic oscillations of transverse momentum distributions,''
\newblock arXiv:1403.3508, 2014,
\newblock \url{https://arxiv.org/abs/1403.3508}.

\bibitem{Wilk2015}
G.~Wilk and Z.~W{\l}odarczyk,
\newblock ``Tsallis Distribution Decorated With Log-Periodic Oscillation,''
\newblock arXiv:1501.02608, 2015,
\newblock \url{https://arxiv.org/abs/1501.02608}.

\bibitem{Giudice2018}
G.~F. Giudice, Y.~Kats, M.~McCullough, R.~Torre, and A.~Urbano,
\newblock ``Clockwork / Linear Dilaton: Structure and Phenomenology,''
\newblock \emph{JHEP} 06 (2018) 009,
\newblock arXiv:1711.08437,
\newblock \url{https://arxiv.org/abs/1711.08437}.

\end{thebibliography}

\end{document}
